\documentclass[11pt]{article}
\usepackage[a4paper,margin=27mm]{geometry}
\usepackage{amsmath,amssymb,hyperref}
\begin{document}
\title{Finite denominator searches for an optimal degree-four rule}
\author{Research note, 30 September 2026}
\date{}\maketitle

The bound in \path{optimized_two_adic_filter.tex} gives $R_4\ge16$,
where repeated rational nodes are allowed. Its equality-case corollary
shows that a putative $16$-node rule has only odd denominators.
If $D$ is a common denominator, the second and fourth moments force
$3\mid D$ and $5\mid D$. Thus $D$ is an odd multiple of $15$.

We exhaustively searched the grids $x_i\in\{0,1/D,\ldots,1\}$ for
$D=15,45,75,105$, finding no rule. This is a finite-grid result only;
it does not improve the universal bound $R_4\ge16$.

Set $z_i=D(2x_i-1)$. These are odd integers in $[-D,D]$, and the
required integer equations are
\[
 \sum z_i=\sum z_i^3=0,\qquad
 \sum z_i^2=\frac{16D^2}{3},\qquad
 \sum z_i^4=\frac{16D^4}{5}.
\]
The program enumerates sorted multisets of sixteen positive magnitudes,
split into their smaller and larger halves of eight. It matches their
second and fourth power sums exactly. For each resulting multiset,
a second meet-in-the-middle search enumerates all signs, matching first
and third power sums. Repetitions and endpoints are included.
The pruning inequalities use only the order of the two halves and
lower or upper bounds for their remaining magnitudes.

\[
\begin{array}{r|r|r|r}
 D&\text{stored smaller halves}&\text{visited larger halves}
  &\text{exact even-moment candidates}\\\hline
45&391619&3021351&2235\\
75&13189332&113553961&30120\\
105&149835139&1358156584&394429
\end{array}
\]
Every candidate failed the exhaustive sign assignment check.
The $D=15$ grid is contained in each of these grids.
The separate, less optimized program
\path{search_degree4_magnitudes.cpp} independently reproduced all
$2235$ candidates at $D=45$ and the same negative result.

The main implementation is \path{search_degree4_ordered_halves.cpp};
the machine-readable report is \path{degree4_grid_search.json}.
No floating-point calculation enters either exact search.

\end{document}
